% ===================================================================
%  Семинар 1. Содержательная часть — единый источник для обоих файлов.
%
%  Состав задач взят из «Семинары.docx»:
%      2.1, 2.3, 2.5, 2.6, 2.12, 2.2, 2.8
%      Источник: Сборник задач к начальному курсу эконометрики (РАНХиГС),
%                гл. 2 «Модель парной регрессии».
%      + отдельная задача, источник:
%        http://www.lib.uniyar.ac.ru/edocs/iuni/20040802.pdf
%
%  Порядок задач сохранён ровно тот, что указан в «Семинары.docx».
%  Раздел «Основные формулы» и всё внутри \begin{solution}...\end{solution}
%  попадает только в seminar_01_solutions.pdf.
%
%  Условия задач — дословно по сборнику. Решения опираются на формулы
%  раздела 1 и ссылаются на них, не повторяя выводов.
% ===================================================================

\seminartitle{Семинар 1}{Модель парной регрессии: МНК, линеаризация,\\ проверка гипотез}

% ===================================================================
% Теория печатается только в файле с решениями.
\ifsolutions
\section{Основные формулы}

% -------------------------------------------------------------------
\subsection{Модель и предпосылки}\label{sec:model}

Модель парной регрессии:
\begin{equation}
  Y_t = \alpha + \beta X_t + \varepsilon_t, \qquad t = 1,\dots,n,
  \label{eq:model}
\end{equation}
где $Y_t$ — зависимая переменная, $X_t$ — объясняющая переменная (регрессор),
$\varepsilon_t$ — случайная ошибка. Параметры $\alpha$ и $\beta$ —
\emph{истинные} значения коэффициентов, то есть неизвестные константы, которые
оцениваются по выборке. Прямую $\alpha + \beta X$ называют истинной линией
регрессии. Ошибки $\varepsilon_t = Y_t - \alpha - \beta X_t$ — отклонения
наблюдений от неё — не наблюдаются.

\textbf{Предпосылки классической модели} (Магнус, с.~39):
\begin{enumerate}[label=\arabic*), leftmargin=2.2em, itemsep=2pt]
  \item выполнена спецификация \eqref{eq:model};
  \item $X_t$ — детерминированные величины, не все равные между собой;
  \item[3a)] $\E\varepsilon_t = 0$ и $\D(\varepsilon_t) = \sigma^2$ при всех $t$;
  \item[3b)] $\Cov(\varepsilon_t, \varepsilon_s) = 0$ при $t \neq s$;
  \item[3c)] $\varepsilon_t \sim N(0, \sigma^2)$.
\end{enumerate}
Свойства оценок в п.~\ref{sec:props}--\ref{sec:s2} выводятся из предпосылок
1--3b. Нормальность 3c нужна только для распределений статистик в
п.~\ref{sec:tests}.

\textbf{Обозначения.} Суммы берутся по $t = 1,\dots,n$.
\begin{itemize}[leftmargin=2.2em, itemsep=2pt]
  \item $\ah$, $\bh$ — оценки параметров, $\yh_t = \ah + \bh X_t$ — подогнанные
        значения, $e_t = Y_t - \yh_t$ — остатки. В отличие от ошибок, остатки
        отсчитываются от \emph{оценённой} прямой и вычисляются по данным.
  \item $\overline{X} = \tfrac{1}{n}\sum X_t$; так же обозначаются средние
        $\overline{Y}$ и $\overline{\varepsilon}$.
  \item $x_t = X_t - \overline{X}$, $y_t = Y_t - \overline{Y}$,
        $\widehat{y}_t = \yh_t - \overline{Y}$ — отклонения от средних.
\end{itemize}

\textbf{Свойства математического ожидания и дисперсии.} Пусть
$Z, Z_1, \dots, Z_n$ — случайные величины, $c, c_1, \dots, c_n$ — константы.
Тогда
\begin{gather}
  \E\Bigl(\sum c_tZ_t\Bigr) = \sum c_t\,\E Z_t,
  \label{eq:E-lin} \\
  \D(Z) = \E\bigl[(Z - \E Z)^2\bigr] = \E\bigl[Z^2\bigr] - \bigl(\E Z\bigr)^2,
  \label{eq:D-def} \\
  \D(c + Z) = \D(Z),
  \qquad
  \D(cZ) = c^2\,\D(Z),
  \qquad
  \D\Bigl(\sum c_tZ_t\Bigr) = \sum c_t^2\,\D(Z_t),
  \label{eq:D-props}
\end{gather}
где последнее равенство верно для попарно некоррелированных $Z_t$.

% -------------------------------------------------------------------
\subsection{Метод наименьших квадратов}\label{sec:ols}

МНК-оценки $\ah$, $\bh$ минимизируют сумму квадратов отклонений наблюдений от
прямой:
\[
  S(a,b) = \sum\bigl(Y_t - a - bX_t\bigr)^2 \;\longrightarrow\; \min_{a,\,b}.
\]
Условия первого порядка
\[
  \frac{\partial S}{\partial a} = -2\sum\bigl(Y_t - a - bX_t\bigr) = 0,
  \qquad
  \frac{\partial S}{\partial b} = -2\sum\bigl(Y_t - a - bX_t\bigr)X_t = 0
\]
в точке $(\ah, \bh)$ записываются через остатки и называются
\emph{нормальными уравнениями}:
\begin{equation}
  \sum e_t = 0, \qquad \sum e_tX_t = 0 .
  \label{eq:normal}
\end{equation}
В развёрнутом виде это система
\[
  \begin{cases}
    n\,\ah + \bh\sum X_t = \sum Y_t, \\[2pt]
    \ah\sum X_t + \bh\sum X_t^2 = \sum X_tY_t .
  \end{cases}
\]
Первое уравнение, делённое на $n$, даёт $\ah = \overline{Y} - \bh\overline{X}$.
Подставив его во второе и учитывая $\sum X_t = n\overline{X}$, получаем
\[
  \bh\Bigl(\sum X_t^2 - n\overline{X}^2\Bigr)
  = \sum X_tY_t - n\overline{X}\,\overline{Y}.
\]
Выражения в скобках — суммы в отклонениях:
\[
  \sum x_t^2 = \sum\bigl(X_t - \overline{X}\bigr)^2
  = \sum X_t^2 - 2\overline{X}\sum X_t + n\overline{X}^2
  = \sum X_t^2 - n\overline{X}^2,
\]
и аналогично $\sum x_ty_t = \sum X_tY_t - n\overline{X}\,\overline{Y}$,
$\sum y_t^2 = \sum Y_t^2 - n\overline{Y}^2$. Так как
$n\overline{X}^2 = \bigl(\sum X_t\bigr)^2/n$, через исходные суммы
\begin{equation}
\begin{gathered}
  \sum x_t^2 = \sum X_t^2 - \frac{\bigl(\sum X_t\bigr)^2}{n},
  \qquad
  \sum y_t^2 = \sum Y_t^2 - \frac{\bigl(\sum Y_t\bigr)^2}{n},
  \\[4pt]
  \sum x_ty_t = \sum X_tY_t - \frac{\sum X_t\sum Y_t}{n}.
\end{gathered}
\label{eq:dev-sums}
\end{equation}
Отсюда МНК-оценки:
\begin{equation}
  \bh = \frac{\sum x_ty_t}{\sum x_t^2}
      = \frac{n\sum X_tY_t - \sum X_t\sum Y_t}{n\sum X_t^2 - \bigl(\sum X_t\bigr)^2},
  \qquad
  \ah = \overline{Y} - \bh\,\overline{X}.
  \label{eq:ols}
\end{equation}
Найденная точка — минимум: матрица вторых производных
\[
  \begin{pmatrix}
    \partial^2S/\partial a^2 & \partial^2S/\partial a\,\partial b \\
    \partial^2S/\partial a\,\partial b & \partial^2S/\partial b^2
  \end{pmatrix}
  = 2\begin{pmatrix} n & \sum X_t \\ \sum X_t & \sum X_t^2 \end{pmatrix}
\]
положительно определена: угловой минор $2n > 0$, определитель
$4\bigl(n\sum X_t^2 - (\sum X_t)^2\bigr) = 4n\sum x_t^2 > 0$ по предпосылке~2.

\textbf{Алгебраические свойства МНК-оценок:}
\begin{equation}
  \sum x_t = 0,
  \qquad
  \widehat{y}_t = \bh x_t,
  \qquad
  e_t = y_t - \bh x_t,
  \qquad
  \sum x_ty_t = \bh\sum x_t^2 .
  \label{eq:ols-props}
\end{equation}
Действительно, $\sum x_t = \sum X_t - n\overline{X} = 0$. Прямая проходит через
точку средних, $\overline{Y} = \ah + \bh\overline{X}$, поэтому
$\widehat{y}_t = (\ah + \bh X_t) - (\ah + \bh\overline{X}) = \bh x_t$ и
$e_t = (Y_t - \overline{Y}) - (\yh_t - \overline{Y}) = y_t - \widehat{y}_t$.
Последнее равенство — это \eqref{eq:ols}.

\begin{remark}
Вывод не зависит от того, какая из переменных объясняющая: в обратной регрессии
$X$ на $Y$ наклон равен $\widehat{\gamma} = \sum x_ty_t \big/ \sum y_t^2$.
\end{remark}

% -------------------------------------------------------------------
\subsection{Разложение суммы квадратов и коэффициент детерминации}\label{sec:decomp}

В регрессии со свободным членом полный разброс зависимой переменной
раскладывается на объяснённую и остаточную части:
\begin{equation}
  \underbrace{\sum\bigl(Y_t - \overline{Y}\bigr)^2}_{\TSS}
  =
  \underbrace{\sum\bigl(\yh_t - \overline{Y}\bigr)^2}_{\RSS}
  +
  \underbrace{\sum\bigl(Y_t - \yh_t\bigr)^2}_{\ESS}.
  \label{eq:decomp}
\end{equation}
Здесь $\TSS$ (total sum of squares) — полная сумма квадратов, $\RSS$
(regression sum of squares) — объяснённая регрессией, $\ESS$ (error sum of
squares) — остаточная, $\ESS = \sum e_t^2$.

\emph{Доказательство.} По \eqref{eq:ols-props} все три суммы выражаются через
отклонения:
\begin{equation}
\begin{gathered}
  \TSS = \sum y_t^2,
  \qquad
  \RSS = \sum\widehat{y}_t^2 = \bh^2\sum x_t^2,
  \\[4pt]
  \ESS = \sum e_t^2 = \sum y_t^2 - \bh^2\sum x_t^2 ,
\end{gathered}
\label{eq:ss-dev}
\end{equation}
где последнее равенство получается раскрытием квадрата:
\[
\begin{aligned}
  \sum\bigl(y_t - \bh x_t\bigr)^2
  &= \sum y_t^2 - 2\bh\sum x_ty_t + \bh^2\sum x_t^2 \\
  &= \sum y_t^2 - 2\bh^2\sum x_t^2 + \bh^2\sum x_t^2
   = \sum y_t^2 - \bh^2\sum x_t^2 .
\end{aligned}
\]
Сумма $\RSS$ и $\ESS$ из \eqref{eq:ss-dev} равна $\TSS$. Свойства
\eqref{eq:ols-props}, на которые опирается доказательство, следуют из
\emph{обоих} нормальных уравнений. В регрессии без свободного члена уравнения
$\sum e_t = 0$ нет, и тождество \eqref{eq:decomp}, вообще говоря, не выполняется
(задача~2.8). $\blacksquare$

\begin{remark}
Обозначения $\RSS$ и $\ESS$ в литературе не унифицированы. Мы следуем сборнику
РАНХиГС: $\RSS$ — \emph{объяснённая} (regression), $\ESS$ — \emph{остаточная}
(error) сумма квадратов. В англоязычных учебниках (Greene, Wooldridge,
Stock--Watson) те же буквы часто означают обратное: $\ESS$ — explained,
$\RSS$ — residual. Ориентируйтесь не на аббревиатуру, а на формулу.
\end{remark}

\textbf{Коэффициент детерминации} — доля полного разброса, объяснённая
регрессией. Так как $\RSS, \ESS \ge 0$,
\begin{equation}
  R^2 = \frac{\RSS}{\TSS} = 1 - \frac{\ESS}{\TSS}, \qquad 0 \le R^2 \le 1 .
  \label{eq:r2}
\end{equation}
Подставив \eqref{eq:ss-dev} и $\bh$ из \eqref{eq:ols}, получаем
\begin{equation}
  R^2 = \frac{\bh^2\sum x_t^2}{\sum y_t^2}
      = \Bigl(\frac{\sum x_ty_t}{\sum x_t^2}\Bigr)^{2}\frac{\sum x_t^2}{\sum y_t^2}
      = \frac{\bigl(\sum x_ty_t\bigr)^2}{\sum x_t^2\sum y_t^2}.
  \label{eq:r2-dev}
\end{equation}

\textbf{Выборочные характеристики} (Магнус, приложение МС, п.~6):
\begin{equation}
  s_X^2 = \frac{\sum x_t^2}{n-1},
  \qquad
  s_Y^2 = \frac{\sum y_t^2}{n-1},
  \qquad
  c_{XY} = \frac{\sum x_ty_t}{n-1},
  \qquad
  r_{XY} = \frac{c_{XY}}{s_Xs_Y}.
  \label{eq:sample}
\end{equation}
Множители $1/(n-1)$ в $r_{XY}^2$ сокращаются, и по \eqref{eq:r2-dev}
\[
  r_{XY}^2 = \frac{c_{XY}^2}{s_X^2s_Y^2}
           = \frac{\bigl(\sum x_ty_t\bigr)^2}{\sum x_t^2\sum y_t^2}
           = R^2 .
\]

\textbf{Средняя ошибка аппроксимации} — среднее относительное отклонение
подогнанных значений от наблюдаемых:
\begin{equation}
  \overline{A} = \frac{1}{n}\sum\biggl|\frac{Y_t - \yh_t}{Y_t}\biggr| \cdot 100\,\% .
  \label{eq:approx}
\end{equation}

% -------------------------------------------------------------------
\subsection{Статистические свойства оценки наклона}\label{sec:props}

Вычтем из модели \eqref{eq:model} её среднее по наблюдениям,
$\overline{Y} = \alpha + \beta\overline{X} + \overline{\varepsilon}$:
\begin{equation}
  y_t = \bigl(\alpha + \beta X_t + \varepsilon_t\bigr)
        - \bigl(\alpha + \beta\overline{X} + \overline{\varepsilon}\bigr)
      = \beta x_t + \varepsilon_t - \overline{\varepsilon}.
  \label{eq:y-dev}
\end{equation}
Подставив это в \eqref{eq:ols} и учитывая $\sum x_t = 0$, выразим оценку через
ошибки:
\begin{equation}
\begin{aligned}
  \bh &= \frac{\sum x_ty_t}{\sum x_t^2}
       = \frac{\sum x_t\bigl(\beta x_t + \varepsilon_t - \overline{\varepsilon}\bigr)}{\sum x_t^2}
       = \frac{\beta\sum x_t^2 + \sum x_t\varepsilon_t
               - \overline{\varepsilon}\sum x_t}{\sum x_t^2} \\[4pt]
      &= \beta + \frac{\sum x_t\varepsilon_t}{\sum x_t^2}.
\end{aligned}
\label{eq:bhat-eps}
\end{equation}
Веса $x_t/\sum x_t^2$ неслучайны (предпосылка~2). Поэтому по \eqref{eq:E-lin} и
предпосылке~3a оценка несмещённая:
\begin{equation}
  \E\bh = \beta + \frac{\sum x_t\,\E\varepsilon_t}{\sum x_t^2} = \beta,
  \label{eq:unbiased}
\end{equation}
а по \eqref{eq:D-props} и предпосылкам 3a, 3b её дисперсия равна
\begin{equation}
\begin{aligned}
  \D(\bh) &= \D\Bigl(\beta + \frac{\sum x_t\varepsilon_t}{\sum x_t^2}\Bigr)
           = \D\Bigl(\frac{\sum x_t\varepsilon_t}{\sum x_t^2}\Bigr)
           = \frac{1}{\bigl(\sum x_t^2\bigr)^2}\sum x_t^2\,\D(\varepsilon_t) \\[4pt]
          &= \frac{\sigma^2\sum x_t^2}{\bigl(\sum x_t^2\bigr)^2}
           = \frac{\sigma^2}{\sum x_t^2}.
\end{aligned}
\label{eq:var-b}
\end{equation}
Точность оценки наклона тем выше, чем больше разброс регрессора $\sum x_t^2$.

% -------------------------------------------------------------------
\subsection{Оценка дисперсии ошибок}\label{sec:s2}

Дисперсия \eqref{eq:var-b} содержит неизвестный параметр $\sigma^2$. Будь ошибки
$\varepsilon_t$ наблюдаемы, его несмещённой оценкой была бы
$\frac{1}{n}\sum\varepsilon_t^2$. Остатки же систематически меньше ошибок:
МНК-прямая минимизирует сумму квадратов среди всех прямых и потому даёт не
большую сумму, чем истинная прямая $\alpha + \beta X$:
\[
  \sum e_t^2 = \min_{a,b}\sum\bigl(Y_t - a - bX_t\bigr)^2
  \le \sum\bigl(Y_t - \alpha - \beta X_t\bigr)^2 = \sum\varepsilon_t^2 .
\]
Поэтому делитель $n$ дал бы оценку, смещённую вниз. Найдём смещение точно.

\textbf{Лемма.} При предпосылках 3a, 3b
\begin{equation}
  \E\sum\bigl(\varepsilon_t - \overline{\varepsilon}\bigr)^2 = (n-1)\sigma^2 .
  \label{eq:lemma}
\end{equation}

\emph{Доказательство.} Раскроем квадрат, учитывая
$\sum\varepsilon_t = n\overline{\varepsilon}$:
\[
  \sum\bigl(\varepsilon_t - \overline{\varepsilon}\bigr)^2
  = \sum\varepsilon_t^2 - 2\overline{\varepsilon}\sum\varepsilon_t + n\overline{\varepsilon}^2
  = \sum\varepsilon_t^2 - 2n\overline{\varepsilon}^2 + n\overline{\varepsilon}^2
  = \sum\varepsilon_t^2 - n\overline{\varepsilon}^2 .
\]
Так как $\E\varepsilon_t = \E\overline{\varepsilon} = 0$, по \eqref{eq:D-def} и
\eqref{eq:D-props}
\[
  \E\bigl[\varepsilon_t^2\bigr] = \D(\varepsilon_t) = \sigma^2,
  \qquad
  \E\bigl[\overline{\varepsilon}^2\bigr]
  = \D\Bigl(\frac{1}{n}\sum\varepsilon_t\Bigr)
  = \frac{1}{n^2}\sum\D(\varepsilon_t)
  = \frac{1}{n^2}\cdot n\sigma^2 = \frac{\sigma^2}{n},
\]
откуда
\[
  \E\sum\bigl(\varepsilon_t - \overline{\varepsilon}\bigr)^2
  = \sum\E\bigl[\varepsilon_t^2\bigr] - n\,\E\bigl[\overline{\varepsilon}^2\bigr]
  = n\sigma^2 - n\cdot\frac{\sigma^2}{n} = (n-1)\sigma^2 .
  \qquad\blacksquare
\]

\textbf{Теорема.} При предпосылках 1--3b
\[
  \E\sum e_t^2 = (n-2)\sigma^2,
\]
поэтому несмещёнными оценками $\sigma^2$ и $\D(\bh)$ служат
\begin{equation}
  s^2 = \frac{1}{n-2}\sum e_t^2,
  \qquad
  s_{\bh}^2 = \frac{s^2}{\sum x_t^2}.
  \label{eq:s2}
\end{equation}
Величина $s_{\bh}$ называется стандартной ошибкой оценки $\bh$.

\emph{Доказательство.} Возьмём математическое ожидание от $\ESS$ в форме
\eqref{eq:ss-dev}; $x_t$ неслучайны:
\[
  \E\sum e_t^2 = \E\sum y_t^2 - \E\bigl[\bh^2\bigr]\sum x_t^2 .
\]
Первое слагаемое найдём по \eqref{eq:y-dev}. Перекрёстный член равен нулю, так
как $\E(\varepsilon_t - \overline{\varepsilon}) = 0$, а последний член даёт
лемма \eqref{eq:lemma}:
\[
\begin{aligned}
  \E\sum y_t^2
  &= \E\sum\bigl(\beta x_t + \varepsilon_t - \overline{\varepsilon}\bigr)^2
   = \beta^2\sum x_t^2
     + 2\beta\sum x_t\,\E\bigl(\varepsilon_t - \overline{\varepsilon}\bigr)
     + \E\sum\bigl(\varepsilon_t - \overline{\varepsilon}\bigr)^2 \\[4pt]
  &= \beta^2\sum x_t^2 + (n-1)\sigma^2 .
\end{aligned}
\]
Второе слагаемое — по \eqref{eq:D-def}, \eqref{eq:unbiased} и \eqref{eq:var-b}:
\[
  \E\bigl[\bh^2\bigr]\sum x_t^2
  = \Bigl(\D(\bh) + \bigl(\E\bh\bigr)^2\Bigr)\sum x_t^2
  = \Bigl(\frac{\sigma^2}{\sum x_t^2} + \beta^2\Bigr)\sum x_t^2
  = \sigma^2 + \beta^2\sum x_t^2 .
\]
Следовательно,
\[
  \E\sum e_t^2 = \beta^2\sum x_t^2 + (n-1)\sigma^2 - \sigma^2 - \beta^2\sum x_t^2
  = (n-2)\sigma^2,
\]
откуда $\E\bigl[s^2\bigr] = \sigma^2$ и
$\E\bigl[s_{\bh}^2\bigr] = \sigma^2/\sum x_t^2 = \D(\bh)$. $\blacksquare$

Число $n-2$ — это число \emph{степеней свободы} остатков: нормальные уравнения
\eqref{eq:normal} накладывают на $n$ остатков две линейные связи, и свободно
меняться могут лишь $n-2$ из них. Общее правило: делитель равен $n$ минус число
оценённых коэффициентов. Для выборочной дисперсии оценивается одно среднее, и,
как показывает лемма, делитель равен $n-1$. В регрессии с $k$ регрессорами и
свободным членом он равен $n-k-1$.

\begin{remark}
Проверка формулы для $s^2$: при $n=2$ прямая проходит ровно через две точки, все
остатки равны нулю, и делитель $n$ дал бы $s^2 = 0$. Но две точки не несут
информации о $\sigma^2$, и формула с $n-2$ честно даёт $0/0$.
\end{remark}

% -------------------------------------------------------------------
\subsection{Проверка гипотез}\label{sec:tests}

\textbf{Распределения Стьюдента и Фишера.} Если $Z \sim N(0,1)$ и
$V \sim \chi^2(k)$ независимы, то
\begin{equation}
  \frac{Z}{\sqrt{V/k}} \sim t(k),
  \qquad
  \frac{Z^2}{V/k} \sim F(1,k).
  \label{eq:t-def}
\end{equation}

\textbf{Распределение $t$-статистики.} При предпосылке 3c оценка
\eqref{eq:bhat-eps} — линейная комбинация нормальных ошибок, поэтому с учётом
\eqref{eq:unbiased} и \eqref{eq:var-b}
\[
  \bh \sim N\Bigl(\beta,\ \frac{\sigma^2}{\sum x_t^2}\Bigr),
  \qquad
  Z = \frac{\bh - \beta}{\sigma\big/\sqrt{\sum x_t^2}} \sim N(0,1).
\]
Кроме того, можно показать, что $V = (n-2)s^2/\sigma^2 \sim \chi^2(n-2)$ и не
зависит от $\bh$. Подставив $Z$ и $V$ в \eqref{eq:t-def} с $k = n-2$, получаем
статистику, в которой неизвестное $\sigma$ сокращается:
\begin{equation}
  \frac{Z}{\sqrt{V/(n-2)}}
  = \frac{(\bh - \beta)\sqrt{\sum x_t^2}\big/\sigma}{s/\sigma}
  = \frac{\bh - \beta}{s\big/\sqrt{\sum x_t^2}}
  = \frac{\bh - \beta}{s_{\bh}} \sim t(n-2).
  \label{eq:t-stat}
\end{equation}
Замена $\sigma$ на $s$ добавляет случайность в знаменатель, поэтому
распределение — Стьюдента, а не нормальное, с тем же числом степеней свободы
$n-2$, что и у $s^2$.

\textbf{$t$-тест.} Гипотеза $H_0\colon \beta = \beta_0$ проверяется статистикой
\begin{equation}
  t = \frac{\bh - \beta_0}{s_{\bh}},
  \label{eq:ttest}
\end{equation}
которая при $H_0$ по \eqref{eq:t-stat} имеет распределение $t(n-2)$. Гипотеза
отвергается, если $|t| > t_{\text{кр}}$; на 5\,\%-ном уровне значимости
$t_{\text{кр}} = t_{0{,}025}(n-2)$.

\textbf{$F$-тест значимости регрессии.} Гипотеза $H_0\colon \beta = 0$
проверяется статистикой
\begin{equation}
  F = \frac{\RSS}{\ESS/(n-2)} = \frac{R^2}{1-R^2}\,(n-2),
  \label{eq:ftest}
\end{equation}
где второе равенство получено делением числителя и знаменателя на $\TSS$. По
\eqref{eq:ss-dev} и \eqref{eq:s2}
\[
  F = \frac{\bh^2\sum x_t^2}{s^2} = \Bigl(\frac{\bh}{s_{\bh}}\Bigr)^2 = t^2,
\]
где $t$ — статистика \eqref{eq:ttest} при $\beta_0 = 0$. Поэтому при $H_0$ по
\eqref{eq:t-def} $F \sim F(1, n-2)$, и в парной регрессии $F$-тест равносилен
$t$-тесту гипотезы $\beta = 0$. Гипотеза отвергается, если
$F > F_{\text{кр}}(1, n-2)$. В регрессии с $m$ регрессорами
\[
  F = \frac{R^2}{1-R^2}\cdot\frac{n-m-1}{m} \sim F(m,\, n-m-1).
\]

% -------------------------------------------------------------------
\subsection{Линеаризация нелинейных моделей}\label{sec:linear}

Формулы п.~\ref{sec:ols}--\ref{sec:tests} применимы к любой модели, которую
заменой переменных и параметров можно привести к виду \eqref{eq:model}, то есть
к модели, \emph{линейной по параметрам} и с \emph{аддитивной} ошибкой.
Нелинейность по переменным этому не мешает. Модели с мультипликативной ошибкой
$\varepsilon_t > 0$ линеаризуются логарифмированием:
\begin{center}
\begin{tabular}{lll}
\toprule
Модель & После логарифмирования & Регрессия \\
\midrule
степенная $Y = aX^{b}\varepsilon$
  & $\ln Y = \ln a + b\ln X + \ln\varepsilon$ & $\ln Y$ на $\ln X$ \\
показательная $Y = a\,b^{X}\varepsilon$
  & $\ln Y = \ln a + X\ln b + \ln\varepsilon$ & $\ln Y$ на $X$ \\
экспоненциальная $Y = a\,e^{bX}\varepsilon$
  & $\ln Y = \ln a + bX + \ln\varepsilon$ & $\ln Y$ на $X$ \\
\bottomrule
\end{tabular}
\end{center}
Исходные параметры находят обратным преобразованием оценок: если $C$ — оценка
свободного члена, то $a = e^{C}$, а при десятичных логарифмах $a = 10^{C}$.
Показатель $b$ степенной модели — эластичность: рост $X$ на 1\,\% меняет $Y$ в
среднем на $b\,\%$.

Если ошибка входит аддитивно, $Y = f(X) + \varepsilon$, логарифм суммы не
расщепляется, и модель не линеаризуется, даже когда сама функция $f$ приводится
к линейному виду.

Качество нелинейной модели оценивают по исходной переменной: подогнанные
значения $\yh_t$ получают обратным преобразованием, а индекс детерминации
$\rho^2 = 1 - \ESS/\TSS$ с $\ESS = \sum\bigl(Y_t - \yh_t\bigr)^2$ и ошибку
$\overline{A}$ считают по \eqref{eq:r2} и \eqref{eq:approx}.
\fi

% ===================================================================
\section{Задачи}

% -------------------------------------------------------------------
\problem{2.1}

Наблюдения 16 пар $(X,Y)$ дали следующие результаты:
\[
  \textstyle\sum Y^2 = 526, \qquad \sum X^2 = 657, \qquad \sum XY = 492,
\]
\[
  \textstyle\sum Y = 64, \qquad \sum X = 96 .
\]
Оцените регрессию $Y_t = \alpha + \beta X_t + \varepsilon_t$ и проверьте
гипотезу, что коэффициент $\beta$ равен $1{,}0$.

\begin{solution}
\textbf{1) Суммы в отклонениях.} Здесь $n = 16$, $\overline{X} = 96/16 = 6$,
$\overline{Y} = 64/16 = 4$. По \eqref{eq:dev-sums}
\[
\begin{aligned}
  \sum x_t^2  &= 657 - \frac{96^2}{16} = 657 - 576 = 81, \\[2pt]
  \sum x_ty_t &= 492 - \frac{96\cdot 64}{16} = 492 - 384 = 108, \\[2pt]
  \sum y_t^2  &= 526 - \frac{64^2}{16} = 526 - 256 = 270 .
\end{aligned}
\]

\medskip
\textbf{2) Оценки параметров.} По \eqref{eq:ols}
\[
  \bh = \frac{108}{81} = \frac{4}{3} \approx 1{,}33,
  \qquad
  \ah = 4 - \frac{4}{3}\cdot 6 = -4 .
\]

\medskip
\textbf{3) Оценка дисперсии ошибок.} По \eqref{eq:ss-dev} и \eqref{eq:s2}
\[
  \sum e_t^2 = 270 - \Bigl(\frac{4}{3}\Bigr)^2\cdot 81 = 270 - 144 = 126,
  \qquad
  s^2 = \frac{126}{14} = 9,
\]
\[
  s_{\bh}^2 = \frac{9}{81} = \frac{1}{9},
  \qquad
  s_{\bh} = \frac{1}{3}.
\]

\medskip
\textbf{4) Проверка гипотезы} $H_0\colon \beta = 1$. По \eqref{eq:ttest}
\[
  t = \frac{4/3 - 1}{1/3} = 1 .
\]
Критическое значение $t_{0{,}025}(14) = 2{,}145$. Так как $|t| = 1 < 2{,}145$,
гипотеза $H_0$ \textbf{не отвергается} на 5\,\%-ном уровне значимости.
\end{solution}

% -------------------------------------------------------------------
\problem{2.3}

Пусть $\bh$ — оценка коэффициента наклона в регрессии $Y$ на $X$,
а $\widehat{\gamma}$ — оценка коэффициента наклона в регрессии $X$ на $Y$.
Покажите, что $\bh = 1/\widehat{\gamma}$ тогда и только тогда, когда $R^2 = 1$.

\begin{solution}
По \eqref{eq:ols} и замечанию в п.~\ref{sec:ols}
\[
  \bh = \frac{\sum x_ty_t}{\sum x_t^2},
  \qquad
  \widehat{\gamma} = \frac{\sum x_ty_t}{\sum y_t^2},
\]
откуда по \eqref{eq:r2-dev}
\[
  \bh\,\widehat{\gamma} = \frac{\bigl(\sum x_ty_t\bigr)^2}{\sum x_t^2\sum y_t^2} = R^2 .
\]
Следовательно,
$\bh = 1/\widehat{\gamma} \iff \bh\,\widehat{\gamma} = 1 \iff R^2 = 1$.

\medskip
\textbf{Геометрический смысл.} $R^2 = 1$ означает $\ESS = 0$, то есть все точки
$(X_t, Y_t)$ лежат на одной прямой. Только в этом случае прямые регрессий $Y$ на
$X$ и $X$ на $Y$ совпадают; в остальных случаях $\bh\,\widehat{\gamma} = R^2 < 1$.
\end{solution}

% -------------------------------------------------------------------
\problem{2.5}

Могут ли следующие уравнения быть преобразованы в уравнения, линейные
по параметрам?
\begin{enumerate}[label=\asbuk*), leftmargin=2.2em, itemsep=3pt]
  \item $Y_i = \alpha \cdot \exp(\beta X_i)\cdot \varepsilon_i$;
  \item $Y_i = \alpha \cdot \exp(-\beta X_i) + \varepsilon_i$;
  \item $Y_i = \exp(\alpha + \beta X_i + \varepsilon_i)$;
  \item $Y_i = \alpha \,/\, (\beta - X_i) + \varepsilon_i$.
\end{enumerate}

\begin{solution}
По п.~\ref{sec:linear} модель линеаризуется, если после преобразования она
линейна по параметрам и ошибка в ней аддитивна.

\medskip
\textbf{а)} Да. Логарифмируя (при $\varepsilon_i > 0$), получаем
\[
  \ln Y_i = \ln\alpha + \beta X_i + \ln\varepsilon_i ,
\]
то есть $Z_i = \gamma + \beta X_i + u_i$, где $Z_i = \ln Y_i$,
$\gamma = \ln\alpha$, $u_i = \ln\varepsilon_i$.

\medskip
\textbf{в)} Да: $\ln Y_i = \alpha + \beta X_i + \varepsilon_i$, то есть
$Z_i = \alpha + \beta X_i + \varepsilon_i$, где $Z_i = \ln Y_i$.

\medskip
\textbf{б), г)} Нет: ошибка входит аддитивно. Сами функции линеаризуются,
\[
  \ln\bigl(\alpha e^{-\beta X_i}\bigr) = \ln\alpha - \beta X_i,
  \qquad
  \frac{1}{\alpha/(\beta - X_i)} = \frac{\beta}{\alpha} - \frac{1}{\alpha}X_i,
\]
но после добавления $\varepsilon_i$ эти преобразования не дают модели, линейной
по параметрам с аддитивной ошибкой.
\end{solution}

% -------------------------------------------------------------------
\problem{2.6}

Зависимая переменная в регрессии $Y_t = \alpha + \beta X_t + \varepsilon_t$
разбивается на две компоненты: $Y_t = Y_{1t} + Y_{2t}$. Рассмотрим две регрессии
для компонент:
\[
  Y_{1t} = \alpha_1 + \beta_1 X_t + \varepsilon_{1t},
  \qquad
  Y_{2t} = \alpha_2 + \beta_2 X_t + \varepsilon_{2t}.
\]
Докажите следующие соотношения для МНК-оценок параметров трёх регрессий:
\[
  \ah = \widehat{\alpha}_1 + \widehat{\alpha}_2,
  \qquad
  \bh = \widehat{\beta}_1 + \widehat{\beta}_2 .
\]

\begin{solution}
Так как $\overline{Y} = \overline{Y}_1 + \overline{Y}_2$, отклонения тоже
складываются: $y_t = y_{1t} + y_{2t}$. По \eqref{eq:ols}
\[
  \bh = \frac{\sum x_t\bigl(y_{1t} + y_{2t}\bigr)}{\sum x_t^2}
      = \frac{\sum x_ty_{1t}}{\sum x_t^2} + \frac{\sum x_ty_{2t}}{\sum x_t^2}
      = \widehat{\beta}_1 + \widehat{\beta}_2,
\]
\[
\begin{aligned}
  \ah &= \overline{Y} - \bh\overline{X}
       = \bigl(\overline{Y}_1 + \overline{Y}_2\bigr)
         - \bigl(\widehat{\beta}_1 + \widehat{\beta}_2\bigr)\overline{X} \\[4pt]
      &= \bigl(\overline{Y}_1 - \widehat{\beta}_1\overline{X}\bigr)
       + \bigl(\overline{Y}_2 - \widehat{\beta}_2\overline{X}\bigr)
       = \widehat{\alpha}_1 + \widehat{\alpha}_2 .
\end{aligned}
\]
\end{solution}

% -------------------------------------------------------------------
\problem{2.12}

Рассмотрим модель регрессии на константу
\[
  Y_t = \alpha + \varepsilon_t, \qquad t = 1,\dots,n .
\]
\begin{enumerate}[label=\asbuk*), leftmargin=2.2em, itemsep=3pt]
  \item Найдите оценки метода наименьших квадратов для $\alpha$ и $\sigma^2$.
  \item Найдите дисперсию оценки $\ah$.
  \item Покажите, что статистика $\dfrac{\ah - \alpha}{s_{\ah}}$ имеет
        распределение $t(n-1)$.
  \item Чему равен коэффициент детерминации $R^2$?
\end{enumerate}

\begin{solution}
\textbf{а)} Минимизируем $S(a) = \sum(Y_t - a)^2$:
\[
  S'(a) = -2\sum(Y_t - a) = 0
  \quad\Longrightarrow\quad
  \ah = \frac{1}{n}\sum Y_t = \overline{Y};
\]
это минимум, так как $S''(a) = 2n > 0$. Остатки
\[
  e_t = Y_t - \overline{Y} = (\alpha + \varepsilon_t) - (\alpha + \overline{\varepsilon})
      = \varepsilon_t - \overline{\varepsilon},
\]
поэтому по лемме \eqref{eq:lemma} $\E\sum e_t^2 = (n-1)\sigma^2$, и несмещённая
оценка $\sigma^2$ равна
\[
  s^2 = \frac{1}{n-1}\sum e_t^2 = \frac{1}{n-1}\sum\bigl(Y_t - \overline{Y}\bigr)^2 .
\]

\medskip
\textbf{б)} По \eqref{eq:D-props}, так как $\D(Y_t) = \D(\alpha + \varepsilon_t) = \sigma^2$
и $Y_t$ некоррелированы,
\[
  \D(\ah) = \D\Bigl(\frac{1}{n}\sum Y_t\Bigr) = \frac{1}{n^2}\sum\D(Y_t)
  = \frac{1}{n^2}\cdot n\sigma^2 = \frac{\sigma^2}{n},
  \qquad
  s_{\ah}^2 = \frac{s^2}{n}.
\]

\medskip
\textbf{в)} При предпосылке 3c
\[
  \ah = \alpha + \overline{\varepsilon} \sim N\Bigl(\alpha,\ \frac{\sigma^2}{n}\Bigr),
  \qquad
  Z = \frac{(\ah - \alpha)\sqrt{n}}{\sigma} \sim N(0,1).
\]
Для нормальной выборки $V = (n-1)s^2/\sigma^2 \sim \chi^2(n-1)$ и не зависит от
$\overline{Y}$ (теорема Фишера). По \eqref{eq:t-def} с $k = n-1$
\[
  \frac{Z}{\sqrt{V/(n-1)}} = \frac{(\ah - \alpha)\sqrt{n}/\sigma}{s/\sigma}
  = \frac{\ah - \alpha}{s/\sqrt{n}} = \frac{\ah - \alpha}{s_{\ah}} \sim t(n-1).
\]

\medskip
\textbf{г)} Подогнанные значения постоянны, $\yh_t = \ah = \overline{Y}$, поэтому
$\RSS = \sum\bigl(\yh_t - \overline{Y}\bigr)^2 = 0$ и по \eqref{eq:r2} $R^2 = 0$.
\end{solution}

% -------------------------------------------------------------------
\problem{2.2}

Покажите, что $\bh = r_{XY}\dfrac{s_Y}{s_X}$, где $r_{XY}$ — выборочный
коэффициент корреляции между $X$ и $Y$, а $s_X$, $s_Y$ — стандартные отклонения
$X$ и $Y$ соответственно.

\begin{solution}
По определениям \eqref{eq:sample} и формуле \eqref{eq:ols}
\[
  r_{XY}\frac{s_Y}{s_X}
  = \frac{c_{XY}}{s_Xs_Y}\cdot\frac{s_Y}{s_X}
  = \frac{c_{XY}}{s_X^2}
  = \frac{\sum x_ty_t/(n-1)}{\sum x_t^2/(n-1)}
  = \frac{\sum x_ty_t}{\sum x_t^2}
  = \bh .
\]
\end{solution}

% -------------------------------------------------------------------
\problem{2.8}

Выведите непосредственно формулу для оценки коэффициента наклона в регрессии
\textbf{без свободного члена}, т.\,е. найдите оценку параметра $\beta$ в регрессии
\[
  Y_t = \beta X_t + \varepsilon_t
\]
минимизацией суммы квадратов отклонений $\sum \bigl(Y_t - \yh_t\bigr)^2$.

\begin{solution}
Минимизируем
\[
  S(b) = \sum\bigl(Y_t - bX_t\bigr)^2 .
\]
Условие первого порядка
\[
  S'(b) = -2\sum\bigl(Y_t - bX_t\bigr)X_t = 0
\]
даёт МНК-оценку
\[
  \boxed{\ \bh = \frac{\sum X_tY_t}{\sum X_t^2}\ }
\]
Это минимум, так как $S''(b) = 2\sum X_t^2 > 0$.

\medskip
Свойства оценки выводятся так же, как в п.~\ref{sec:props}, с $X_t$ вместо
$x_t$. Подставим $Y_t = \beta X_t + \varepsilon_t$:
\[
  \bh = \frac{\sum X_t\bigl(\beta X_t + \varepsilon_t\bigr)}{\sum X_t^2}
      = \frac{\beta\sum X_t^2 + \sum X_t\varepsilon_t}{\sum X_t^2}
      = \beta + \frac{\sum X_t\varepsilon_t}{\sum X_t^2}.
\]
Веса $X_t/\sum X_t^2$ неслучайны, поэтому по \eqref{eq:E-lin} и предпосылке~3a
$\E\bh = \beta$, а по \eqref{eq:D-props} и предпосылкам 3a, 3b
\[
\begin{aligned}
  \D(\bh) &= \D\Bigl(\beta + \frac{\sum X_t\varepsilon_t}{\sum X_t^2}\Bigr)
           = \frac{1}{\bigl(\sum X_t^2\bigr)^2}\,\D\Bigl(\sum X_t\varepsilon_t\Bigr)
           = \frac{1}{\bigl(\sum X_t^2\bigr)^2}\sum X_t^2\,\D(\varepsilon_t) \\[4pt]
          &= \frac{\sigma^2\sum X_t^2}{\bigl(\sum X_t^2\bigr)^2}
           = \frac{\sigma^2}{\sum X_t^2}.
\end{aligned}
\]

\medskip
Условие первого порядка здесь — единственное нормальное уравнение
$\sum e_tX_t = 0$. Уравнения $\sum e_t = 0$ нет, так как нет свободного члена,
поэтому доказательство тождества \eqref{eq:decomp} из п.~\ref{sec:decomp} не
проходит, и в общем случае $\TSS \neq \RSS + \ESS$. Это понадобится в
задаче~2.9 на следующем семинаре.
\end{solution}

% ===================================================================
\section{Дополнительная задача}

\ifsolutions   % примечание об источнике — только в решениях
\begin{remark}
\raggedright
Источник отличается от основного сборника:
\href{http://www.lib.uniyar.ac.ru/edocs/iuni/20040802.pdf}%
{lib.uniyar.ac.ru/edocs/iuni/20040802.pdf}, «Часть 2. Решение типовых задач»,
задача 1.
\end{remark}
\fi

\problem{на нелинейные модели}

По семи территориям Уральского региона за 2002 год известны значения двух
признаков:
\begin{itemize}[leftmargin=2.2em, itemsep=1pt, topsep=3pt]
  \item $y$ — расходы на покупку продовольственных товаров в общих расходах, \%;
  \item $x$ — среднедневная заработная плата одного работающего, руб.
\end{itemize}

\begin{center}
\begin{tabular}{clcc}
\toprule
Номер & Район & $y$ & $x$ \\
\midrule
1 & Удмуртская респ.   & 68,8 & 45,1 \\
2 & Свердловская обл.  & 61,2 & 59,0 \\
3 & Башкортостан       & 59,9 & 57,2 \\
4 & Челябинская обл.   & 56,7 & 61,8 \\
5 & Пермская обл.      & 55,0 & 58,8 \\
6 & Курганская обл.    & 54,3 & 47,2 \\
7 & Оренбургская обл.  & 49,3 & 55,2 \\
\bottomrule
\end{tabular}
\end{center}

\noindent\textbf{Задание.}
\begin{enumerate}[label=\arabic*., leftmargin=2.2em, itemsep=3pt]
  \item Для характеристики зависимости $y$ от $x$ рассчитать параметры следующих
        функций: 1)~линейной; 2)~степенной; 3)~показательной.
  \item Оценить каждую модель через коэффициент детерминации $R^2$, среднюю
        ошибку аппроксимации $\overline{A}$ и $F$-критерий Фишера.
\end{enumerate}

\begin{solution}
\textbf{Вспомогательные величины:}
\[
  n = 7, \quad
  \overline{y} = 57{,}89, \quad
  \overline{x} = 54{,}90, \quad
  \overline{yx} = 3166{,}05, \quad
  \sigma_y = 5{,}74, \quad
  \sigma_x = 5{,}86 .
\]
Здесь $\overline{yx} = \frac{1}{n}\sum y_tx_t$, а $\sigma$ — стандартные
отклонения с делителем $n$. Промежуточные величины округлены; оценки рассчитаны
по неокруглённым данным.

\medskip
\textbf{1) Линейная модель} $\widehat{y} = a + bx$. Формула для $b$ — это
\eqref{eq:ols}, в которой числитель и знаменатель поделены на $n$:
\[
  b = \frac{\overline{yx} - \overline{y}\cdot\overline{x}}{\sigma_x^2} \approx -0{,}346,
  \qquad
  a = \overline{y} - b\,\overline{x} \approx 76{,}88,
\]
\[
  \widehat{y} = 76{,}88 - 0{,}346\,x .
\]
С ростом среднедневной заработной платы на 1 руб. доля расходов на покупку
продовольственных товаров снижается в среднем на 0,35 процентного пункта.
По задаче 2.2, п.~\ref{sec:decomp} и \eqref{eq:approx}
\[
  r_{xy} = b\,\frac{\sigma_x}{\sigma_y} \approx -0{,}353,
  \qquad
  R^2 = r_{xy}^2 \approx 0{,}125,
  \qquad
  \overline{A} = 8{,}1\,\% .
\]
Связь умеренная и обратная; вариацией $x$ объясняется 12,5\,\% вариации $y$.

\medskip
\textbf{2) Степенная модель} $\widehat{y} = a\,x^{b}$. Логарифмируя по
п.~\ref{sec:linear}, получаем $\lg y = \lg a + b\lg x$, то есть $Y = C + bX$,
где $Y = \lg y$, $X = \lg x$, $C = \lg a$. По логарифмам данных
($\overline{Y} = 1{,}7605$, $\overline{X} = 1{,}7370$,
$\overline{YX} = 3{,}0572$, $\sigma_X = 0{,}0484$)
\[
  b = \frac{\overline{YX} - \overline{Y}\cdot\overline{X}}{\sigma_X^2} \approx -0{,}298,
  \qquad
  C = \overline{Y} - b\,\overline{X} \approx 2{,}279,
\]
\[
  \widehat{y} = 10^{2{,}279}\,x^{-0{,}298} \approx 190{,}0\,x^{-0{,}298}.
\]
Рост заработной платы на 1\,\% снижает долю расходов на продовольствие в
среднем на 0,298\,\%. По исходной переменной
\[
  \rho^2 = 1 - \frac{\sigma^2_{\text{ост}}}{\sigma_y^2}
         = 1 - \frac{28{,}32}{32{,}92} \approx 0{,}140,
  \qquad
  \overline{A} = 8{,}0\,\% ,
\]
то есть степенная модель описывает связь несколько лучше линейной.

\medskip
\textbf{3) Показательная модель} $\widehat{y} = a\,b^{x}$. Логарифмируя по
п.~\ref{sec:linear}, получаем $\lg y = \lg a + x\lg b$, то есть $Y = C + Bx$,
где $Y = \lg y$, $C = \lg a$, $B = \lg b$. По тем же формулам для пар
$(x_t, Y_t)$
\[
  B \approx -0{,}00232,
  \qquad
  C = \overline{Y} - B\,\overline{x} \approx 1{,}8878,
\]
\[
  \widehat{y} = 10^{1{,}8878}\cdot\bigl(10^{-0{,}00232}\bigr)^{x}
              \approx 77{,}24\cdot 0{,}9947^{\,x}.
\]
Рост заработной платы на 1 руб. уменьшает долю расходов на продовольствие в
среднем в 0,9947 раза, то есть на 0,53\,\%. По исходной переменной
$\rho^2 \approx 0{,}125$, $\overline{A} = 8{,}1\,\%$.

\medskip
\textbf{Сводка и $F$-критерий.} По \eqref{eq:ftest} при $n = 7$ статистика
равна $F = 5R^2/(1-R^2)$, критические значения $F_{\text{кр}}(1,5) = 6{,}61$ на
уровне 5\,\% и $16{,}26$ на уровне 1\,\%.
\begin{center}
\begin{tabular}{lccc}
\toprule
Модель & $R^2$ (или $\rho^2$) & $\overline{A}$ & $F_{\text{факт}}$ \\
\midrule
Линейная       & 0,125 & 8,1\,\% & 0,71 \\
Степенная      & 0,140 & 8,0\,\% & 0,81 \\
Показательная  & 0,125 & 8,1\,\% & 0,71 \\
\bottomrule
\end{tabular}
\end{center}

\noindent Во всех трёх случаях $F_{\text{факт}} < F_{\text{кр}} = 6{,}61$, поэтому
гипотеза $H_0$ о статистической незначимости уравнения \textbf{не отвергается}:
ни одна из моделей не значима на 5\,\%-ном уровне. Формально лучшая из трёх —
степенная (наибольший $\rho^2$, наименьшая $\overline{A}$), но её преимущество
в пределах погрешности. Причина — слабая связь и очень малое число наблюдений
($n = 7$).

\begin{remark}
В исходном пособии в разборе линейной модели написано, что
$F_{\text{факт}} > F_{\text{кр}}$ при $\alpha = 5\,\%$ и «гипотеза $H_0$
не принимается, что говорит о значимости уравнения регрессии в целом».
Это ошибка: там же приведено $F_{\text{факт}} = 0{,}714$, а
$F_{\text{кр}} = 6{,}61$, то есть $F_{\text{факт}} < F_{\text{кр}}$ и уравнение
\emph{незначимо}. Двумя строками ниже пособие само себе противоречит, верно
заключая обратное для $\alpha = 1\,\%$. Хорошее место, чтобы на семинаре
разобрать, в какую сторону читается $F$-тест.
\end{remark}
\end{solution}
